\section{F. Companion statement map and literature record}
\subsection{F.1. Manuscript wording and statement map}
We quote the supplied AI transcription witnesses, not the absent manuscript
scans.  We omit struck-out words and retain inserted words where the resulting
reading is unambiguous.  Cote and scan numbers, rather than the manuscript's
own page numbers, follow the prompt's convention.  A quotation records the
question being addressed; it does not itself certify the mathematical claim.

\par\smallskip\noindent\textbf{A.1, cote 4:4.}
``Nous admettrons que ces relations restent valables dans le cas de la
cohomologie cristalline.''  The assertions are integral injectivity, torsion
identification and torsion-freeness of $E(Y)$.
Theorems~\ref{thm:A-high-degree} and \ref{thm:A-all-degree-ranges}
give the explicit sufficient degree range and two all-degree classes.
Proposition~\ref{prop:A-transpose-u2} types the missing (8.2.1): it is the
torsion Gysin transpose, and after the justified torsion identification its
composite is cup product by $d\,c_1(\mathcal O_X(1))$.
The full small-degree problem is retained in
Question~\ref{open:A-small-degree}.

\par\smallskip\noindent\textbf{A.2, cote 4:8, and the sheet 4:2.}
``Je pense que l'homomorphisme $v$ de (8.5.2) est injectif'' is settled by
Lemma~\ref{lem:A-R1} without a degree bound.
Theorems~\ref{thm:A-exact85} and \ref{thm:A-unit-count} then establish
the stable exact sequence, the dimension formulas and the unit-root count
for every smooth section.  The remark ``Sans doute (*) et (**) sont
transposés l'un de l'autre'' is made precise in
Proposition~\ref{prop:A-star-dual}; the adjoint operators are $F$ and $V$.

\par\smallskip\noindent\textbf{A.3, cote 4:9.}
The retained marginal assertion is ``Or cela est établi au n° 7'', beside
the eventual positivity of $\dim E(Y,\mathcal O_Y)_{\mathrm{ss}}$.
Theorem~\ref{thm:all-X-explicit-positive-rank} supplies an explicit
threshold, and Corollary~\ref{cor:eventual-hodge-parenthesis} settles
the eventual generic Hodge-level parenthesis.  The relative-ordinarity
prediction is refuted by the characteristic-two surface constructed above.
Proposition~\ref{prop:effective-generic-rank} computes the exact generic
rank for specified finite-field input; a structural rank formula for
arbitrary ambient $X$ remains an explicitly retained obligation.

\par\smallskip\noindent\textbf{A.4, cote 4:12.}
The question asks whether the image ``est-elle égale à
$\operatorname{Aut}(V)$ tout entier ?'', with an inserted question on
irreducibility.  Theorems~\ref{thm:pn-odd-full-groups} and
\ref{thm:pn-char2-all-groups} determine every projective-space case.
The latter refutes the unrestricted full-linear prediction in infinitely
many degrees.  Proposition~\ref{prop:exact-general-determinant-image}
determines the geometric and arithmetic determinant images for general
$X$ from the stable Hasse invariant.  The complete nonabelian image for
arbitrary $X$ remains open here.

\par\smallskip\noindent\textbf{A.5, cote 4:12.}
The requested consequence is that ``$\operatorname{card}Y(k')$ soit
congru mod $p$ à un élément de $\underline F_p$ donné à l'avance''.
Theorem~\ref{thm:all-X-residue-density} proves this for every ambient
variety in the question, with its explicit $d_{\rm res}$, in every
nonempty open.  Corollary~\ref{cor:pn-all-p-sharp-residue-threshold}
gives the sharp bound $d=n+2$ in projective space, including $d=4$
for planes.  The low-degree residue cases appear in the plane and
projective-space classifications.

\par\smallskip\noindent\textbf{A.6, cote 4:10--11.}
The expected comparison concerns the ``torsion de cohomologie
$p$-adique de la variété relevée''.
Theorem~\ref{thm:A-special-torsion} gives the exact special-fibre
identification and specialization obstruction.
Theorem~\ref{thm:A-BMS-counterexample} refutes the unrestricted
expectation; Corollary~\ref{cor:A-lifting-small-range} gives the
valid small-ramification implication.  The final ample-section search
in the counterexample is specified by a terminating construction,
not represented as an executed numerical search.

\subsection{Manuscript anchors for Parts B and C}

The quotations below are transcribed from the supplied page files in
the package \path{MPL_8_079-081+9+4+100}. The references use the
\emph{cote:scan} convention, followed by the manuscript page where
the transcription supplies one. Deletions have been omitted and
insertions accepted, except where an explicitly marked uncertainty
matters. These are checks of the supplied transcription, not fresh
readings of the original scans.

\par\smallskip\noindent\textbf{B.1: the supersingular test and the negative marginal answer.}
At 9:13 (manuscript p.12) the proposed test is:
\begin{quote}
``Cas test: soient $E_i$ $(i=1,2)$ deux courbes elliptiques sur $K$,
qui ont bonne réduction essentielle, ${}_pE_1$ isomorphe à
${}_pE_2$, $K$ contient les racines p.èmes de $1$''.
\end{quote}
The same page singles out the case where one curve reduces to
``une courbe à invariant de Hasse nul''. The next page, 9:14
(manuscript p.13), carries the marginal words ``faux [par?] Tate'';
the supplied transcription marks \emph{par} uncertain. The explicit
examples in B-RD1--2 and the ordinary theorem B-RD3 distinguish the
two cases without relying on the earlier proposed Kisin-module
construction.

\par\smallskip\noindent\textbf{B.2: the scope of the small-defect list.}
The hypothesis introducing the list on 9:14 (manuscript p.13) is
essential:
\begin{quote}
``Dans le cas d'une courbe elliptique, et en car. rés. $\ne 2,3$,
prenant pour $K'$ la plus petite extension sur laquelle on a
réduction semi-stable, on trouve que $G$ est un sous-groupe de
$\operatorname{Aut}(A_{s'})$, lequel est cyclique a $2,4$ ou $6$
éléments donc est premier à $p$''.
\end{quote}
This sentence does not assert the same list in residue
characteristic three. B-RD8--10 explain the additional wild
defect~12 and the exact first-level ramification ranges.

\par\smallskip\noindent\textbf{B.3: constant torsion and the trace normalization.}
The question at 9:14 is whether, when
``${}_pA_K$ soit un groupe constant'', the trace is
``égale à $n=\dim A_K$'' on prime-to-$p$ inertia. The rank of the
covariant Dieudonne module is $2\dim A_K$, so this is a normalization
slip; the same page correctly uses the value~$2$ for an elliptic
curve. It explicitly describes the search as
\begin{quote}
``trouver un exemple où le critère de réduction semi-stable de
Raynaud est en défaut pour $\ell=p$''.
\end{quote}
B-RD12--15 consequently distinguish the potentially good elliptic
question from the unrestricted question, where multiplicative
Tate curves are relevant.

\par\smallskip\noindent\textbf{B.4 and Lead B': the precise monodromy condition.}
At 9:11 (manuscript p.10) the definition is:
\begin{quote}
``satisfait à la condition de monodromie s'il admet une suite de
composition par des groupes de B.T., telle que les
$N=\operatorname{Gr}^i(M)$ aient essentiellement bonne réduction''.
\end{quote}
At 9:12 (manuscript p.11) the trace of the resulting inertia action
is called the \emph{caractère}; 9:13 asks if it is determined
``par la seule connaissance de $M(0)={}_pM$''. The comparison is
expressly restricted to elements ``d'ordre premier à $p$''.
B-RD4 shows why the filtration condition cannot silently be replaced
by potential semistability of the whole representation. B-RD17--18
prove finite upper and nonzero lower bounds for this precise
prime-to-$p$ character; the sharp general-height index remains a
separate question.

\par\smallskip\noindent\textbf{C.1: the asserted equivalence and the missing overlap.}
The theorem on 100:3 (manuscript p.2) assumes a finite surjective
$S'\to S$ with $S'$ geometrically unibranch and
``$\operatorname{prof\,et}_T(S)\geq2$'', and concludes:
\begin{quote}
``Alors $F$ est une équivalence des catégories''.
\end{quote}
The reduction step on 100:7 (manuscript p.6) says:
\begin{quote}
``Or $U\amalg S'\to S$ est fid. plat, et on trouve une donnée de
descente sur $A_U\amalg A'$''.
\end{quote}
The page marks this step with question marks. An actual isomorphism
on the completed overlap is an additional requirement; C-RD12
identifies it exactly after a reduced realization has been supplied.

\par\smallskip\noindent\textbf{C.2: nonfinite normalization was a proof obligation.}
At 100:11 (manuscript p.10) the text is explicit:
\begin{quote}
``Le cas où on ne suppose plus $S''$ fini sur $S'$ se traite par la
méthode essentiellement de $[\ ]$ p 75--76, à condition qu'on
prouve que l'isomorphisme sur $U''\simeq U'$ entre
$T_p(A'')\mid U''$ et $M_{S''}\mid U''$ se prolonge à $S''$''.
\end{quote}
The transcription records an inserted ``O.K.'' and a marginal
question mark. An exact printed parallel, and a plausible identification of the missing reference, is
\cite[proof of Corollary~4.2, pp.75--76]{Grothendieck1966Hom}.
Its normal-base theorem does not supply the claimed reduced-base
extension. C-RD11 gives a reduced geometrically unibranch
counterexample with nonfinite normalization.

\par\smallskip\noindent\textbf{C.3: descent of homomorphisms of BT groups.}
At 100:25 (manuscript p.24), Corollary~3 states:
\begin{quote}
``Alors $f$ est un morphisme de descente pour la catégorie des
schémas en groupes $p$-divisibles''.
\end{quote}
Its remaining requirement is ``Il faut prouver que $X$ est
formellement net'', with $X=\underline{\operatorname{Hom}}_S(A,B)$.
C-RD1 proves the descent equalizer in the stated noetherian
setting; C-RD2 records why an unrestricted nonnoetherian rigidity
claim would be false.

\par\smallskip\noindent\textbf{C.4: the local-system lemma and its marginal doubt.}
At 100:17 (manuscript p.16) the proposed lemma reads:
\begin{quote}
``Un `syst local' de schémas abéliens'' sur un schéma [loc. noeth]
provient d'un schéma abélien.''
\end{quote}
The square brackets retain the transcription's uncertain expansion
of the abbreviation. The marginal judgment is unambiguous:
``Douteux, sans doute faux''. C-RD7--10 give the precise division:
effectivity over normal or geometrically unibranch bases, and under
the local finite-cover hypothesis; failure over arbitrary reduced
noetherian bases; and invariance under nilpotent reduction for an
already specified compatible local system.

% F.1 companion ledger, copied from the supplied transcription files.
% Scan numbers, rather than the sometimes shifted manuscript foliation, are used.
\subsection{Manuscript statements settled by Parts D and E}

\begin{remark}[D.1, cote 8, scan 81]\label{rem:DE-quote-D1}
\status{Proved under stated hypotheses}
The page says ``Elle est évidemment nulle si $M_0$ est de rang $\leq1$, mais sans doute pas nulle si $M_0$ de rang $d\geq2$,'' in the universal general-linear case.  Theorem~\ref{thm:DE-split} confirms the obstruction interpretation and corrects this proposed nonsplitting list by the three exceptions $(d,p)=(2,2),(2,3),(3,2)$.  It also gives the requested complete special-linear and symplectic lists.  The distinction between a supplement and a complement is retained.
\end{remark}

\begin{remark}[D.2, cote 8, scans 80--81]\label{rem:DE-quote-D2}
\status{Proved under stated hypotheses}
At scan 81 the transcribed comparison ends in $H^2(S,\operatorname{End}(E))$, while the concluding argument says ``si le rev. $S'$ de $S$ est le rev. universel'' then the group-cohomology edge map ``est injectif.''  These are distinct steps.  Theorem~\ref{thm:DE-coherent-obstruction} constructs the requested smooth projective $S$ with $\xi\ne0$ for every nonsplit pair.  Remark~\ref{rem:DE-universal-negative} constructs full-monodromy smooth projective curves with $\xi=0$, so the universal version of Lead D is false.
\end{remark}

\begin{remark}[E.1 and R11, cote 5, scan 41]\label{rem:DE-quote-E1}
\status{Proved under stated hypotheses}
After the two rules for $\pi$, the question is ``Si on a une application satisfaisant ces conditions, peut on définir les $\pi^{(n)}$'' satisfying the divided-power axioms.  Theorem~\ref{thm:DE-PD-pi} answers yes for every $\mathbb Z_{(p)}$-algebra, including arbitrary $p$-torsion and $p=2$.  The same page asks of $c_n\bmod p$, ``Est-ce tjs $1$?''  The answer is no for $p\geq3$: the exact formula gives $c_n\equiv\prod_i(a_i!)^{-1}\pmod p$ for the base-$p$ digits $a_i$, and already $n=2$ gives $1/2\ne1$.  At $p=2$ all these residues are $1$.
\end{remark}

\begin{remark}[E.2, cote 6, scan 7, section 1.8]\label{rem:DE-quote-E2}
\status{Proved under stated hypotheses}
The exact transcribed sentence reads ``Sauf erreur, on a encore une équivalence de catégories si $S\to R$ est plat et localement de présentation finie,'' followed by the insertion ``(du moins si R loc.noeth., et se bornant aux Modules cohérents)'', then ``mais je n'ai pas écrit'' the demonstration.  The insertion materially narrows the manuscript conjecture.  Proposition~\ref{prop:DE-E2-coherent} proves the coherent characteristic-zero field case.  Proposition~\ref{prop:DE-E2-Fpure} proves the arbitrary quasi-coherent case for $F$-pure schemes locally of finite type over a perfect field of positive characteristic, using Lemma~\ref{lem:DE-crystal-Fdiv}.  The full flat relative assertion requested in Part E.2 is refuted by Theorem~\ref{thm:DE-E2-counterexample}, which gives a non-Noetherian finite free example and a rank-one crystal automorphism invisible to restriction.  Corollary~\ref{cor:DE-E2-Noetherian-counterexample} further gives nonfaithfulness over a regular Noetherian base with arbitrary quasi-coherent coefficients.  The inserted relative coherent conjecture over every locally Noetherian base is proved by Theorem~\ref{thm:DE-E2-coherent-Noetherian}, with ordinary tensor pullback and arbitrary classical tests.
\end{remark}

\begin{remark}[E.3 and R10, cote 7, scan 17]\label{rem:DE-quote-E3}
\status{Proved under stated hypotheses}
The hoped-for comparison with integral de Rham cohomology is followed by ``C'est faux malheureusement'' and the ordinary-abelian calculation
$H^i((X_0/S)_{\mathrm{cris}},\underline O)=0$ for $i>\dim X_0$.  Here the notation $\mathrm{cris}$ means the old infinitesimal site without divided powers.  Theorem~\ref{thm:DE-inf} identifies this cohomology integrally with $H^i_{\et}(X_0,\Zp)\otimes_{\Zp}W$, including torsion.  Proposition~\ref{prop:DE-R10} verifies the displayed vanishing and the preceding mod-$p$ statement.  No identification with modern PD crystalline cohomology is intended.
\end{remark}

\begin{remark}[E.4, cote 7, scan 70]\label{rem:DE-quote-E4}
\status{Proved under stated hypotheses}
The conjecture says ``Pour tout espace analytique $\mathfrak X$, le complexe $\mathcal F_{\mathfrak X}^*$ est une résolution du faisceau $\mathbf C_{\mathfrak X}$.''  Theorem~\ref{thm:DE-intrinsic-nerve} settles the algebraic finite-type characteristic-zero question posed in Part E.4: the intrinsic completed nerve computes infinitesimal de Rham cohomology, including for singular schemes.  For $\mathbb C[t^2,t^3]$, it gives $H^0=\mathbb C$ and $H^i=0$ for $i>0$.  The theorem is not stated as a proof of the broader assertion for every analytic space.
\end{remark}

\begin{remark}[E.5, cote 7, scan 24]\label{rem:DE-quote-E5}
\status{Proved under stated hypotheses}
For a singular complex curve, the page says ``on n'obtient déjà pas pour'' $\dim H^1_{\mathrm{DR}}(X)$ ``le premier nb de Betti comme il faudrait.''  Its proposed ordinary-cusp example does not establish this.  Theorem~\ref{thm:DE-curves} proves the precise correction for reduced affine locally planar complex curves:
\[
 \dim_{\mathbb C}H^1(\Omega_C^\bullet)
 =b_1(C)+\sum_x(\mu_x-\tau_x),\qquad
 H^i(\Omega_C^\bullet)=0\quad(i\geq2).
\]
Thus the extra term vanishes for the ordinary cusp, although it need not vanish for other plane singularities.
\end{remark}

\begin{remark}[E.6, cote 7, scans 42--43]\label{rem:DE-quote-E6}
\status{Proved under stated hypotheses}
The proposed theorem begins ``Il semble que le théorème suivant soit vrai.''  The key assertion at scan 43 is ``Cet homomorphisme ne dépend que de l'homomorphisme $u_0$'' induced modulo $p$; the page notes that the then-available proof was only ``mod torsion, i.e. en tensorisant par $K$.''  Theorem~\ref{thm:DE-integral-MW-independence} confirms the integral assertion, for every prime and in every degree, with the required overconvergent homotopy supplied by Ertl--Sprang.
\end{remark}

\begin{remark}[E.7, cote 8, scans 48--50]\label{rem:DE-quote-E7}
\status{Proved under stated hypotheses}
The notation is $t_G=\pi^{-\alpha}\Omega^\vee$.  At scan 50, after comparing the two full lattices, the page asks ``On peut se demander si ces sous-réseaux sont les mêmes'' and adds ``ce qui impliquerait évidemment'' $\alpha=0$.  Proposition~\ref{prop:DE-LT} proves $\alpha=0$ for every finite extension $K/\Qp$, but refutes equality of the full lattices using $\pi^{p+1}=p$.  The implication on the manuscript page is correct; a converse or equivalence must not be attributed to it.  The dual $\Omega^\vee$ is the $\Zp$-dual, so the inverse different is part of the normalization.
\end{remark}

\begin{remark}[E.8, cote 9, scan 22]\label{rem:DE-quote-E8}
\status{Proved under stated hypotheses}
The assertion is ``cet isom est compatible avec les $p^{-1}$-structures, donc transforme $\sigma_A$ en $\sigma_{A^*}$.''  Theorem~\ref{thm:DE-hasse} proves exact transport of the Hasse invariant under the stated Hodge-determinant and trace identification, for every abelian scheme over an $\Fp$-scheme, with no reducedness or ordinarity assumption.
\end{remark}

\begin{remark}[E.9, cote 9, scan 29]\label{rem:DE-quote-E9}
\status{Proved under stated hypotheses}
Part c) reads ``Il faudrait vérifier que'' $\alpha_{A^*}=\alpha_A^{-1}$, followed by ``[ aux signes (?)\ldots].''  Part d) compares the reduction of $\alpha_A$ with the map ``déduit plus haut de la dualité en cohomologie de DeRham (à vérifier \ldots).''  Theorem~\ref{thm:DE-ordinary-sign} proves exact inverse biduality for standard Poincaré biduality, and computes the de Rham comparison factor
$\varepsilon_g=(-1)^{g(g+1)/2}$ with its required direction/inverse.  In particular the literal lead's simple multiplicativity is false in odd characteristic: a product contributes the additional factor $(-1)^{gh}$.
\end{remark}

\begin{remark}[E.10, cote 7, scans 18--19]\label{rem:DE-quote-E10}
\status{Proved under stated hypotheses}
The heading is ``Tests décisifs'': for smooth affine $X_0$, ``on doit retrouver la définition de W.M.''; for smooth proper liftable $X_0$, ``on doit retrouver (du moins mod torsion) la cohomologie de De Rham''.  The page calls this ``une première'' ``esquisse.''  Proposition~\ref{prop:DE-old-site}, Theorem~\ref{thm:DE-old-Gm}, and Proposition~\ref{prop:DE-old-proper} refute the integral affine test, its rational version, and the proper test, respectively, on the literal flat weak-formal site.  The rational rigid-cohomology claim of Lead E10 is therefore false as well.
\end{remark}

\begin{remark}[E.11, cote 6, scan 14]\label{rem:DE-quote-E11}
\status{Proved under stated hypotheses}
The page concludes ``On voit donc là un principe d'approche transcendante pour l'étude des familles de groupes formels'' in characteristic $p$.  Theorem~\ref{thm:DE-simply-connected-convergent} proves the requested statement for convergent $F$-isocrystals under the specified proper smooth algebraic lift and simply connected complex fibre.  Corollary~\ref{cor:DE-BT-constant} proves the additional isogeny-to-constant assertion for Barsotti--Tate groups in the literal lead.  These hypotheses retain both the specified embedding $W(k)\hookrightarrow\mathbb C$ and the condition on the generic complex fibre.
\end{remark}

% Supplemental F.1 audit. All four supplied proofs were read before this check.
\subsection{The remaining background checks: R7--R10}

\begin{proposition}[R7, scans 9:16--18]\label{prop:DE-R7}
\status{Proved under stated hypotheses}
Let $A_0/k$ be an abelian variety in characteristic $p$.  A base-change-compatible, functorial ordinary $k$-stratification on $H^1_{\dR}(A/\Lambda)$ for every infinitesimal deformation of $A_0$, horizontal for the relative operators
$F:H^{(p)}\to H$ and $V:H\to H^{(p)}$, exists if and only if $A_0$ is ordinary.  In that case it is unique and preserves the canonical splitting
\[
 H=\omega_{A/\Lambda}\oplus\im F.
\]
Thus the Hodge filtration is horizontal, as asserted in the manuscript.
\end{proposition}
\begin{proof}
The formally smooth universal unpolarized deformation space is $\operatorname{Spf}k[[t_1,\ldots,t_{g^2}]]$; this follows from the Barsotti--Tate deformation theorem and Serre--Tate, in \cite[Corollary~4.8(i)]{Illusie1985BT} and \cite[Theorem~1.2.1]{Katz1981ST}.  Pulling a formal stratification back by $(0,t)$ trivializes its underlying module over this complete local ring.  Horizontality then identifies every iterate of $F$ and $V$ with the corresponding operator on the closed fibre.  Their geometric stable ranks cannot change.

For clarity, the ordinary-deformation input can be verified from an opened explicit display, without an unlocated assertion about density.  After extending the residue field to an algebraic closure, split off the \etale part of $A_0[p^\infty]$.  For the remaining formal group of height $h$ and dimension $d$, write Oort's covariant display matrix as
\[
 \begin{pmatrix}A&B\\C&D\end{pmatrix}\in\GL_h(W(k)),\qquad
 F=\begin{pmatrix}A&pB\\C&pD\end{pmatrix}\sigma.
\]
The construction in \cite[\S\S1.9, 1.12--1.14]{Oort2000} replaces the first matrix by
$\left(\begin{smallmatrix}A+TC&B+TD\\C&D\end{smallmatrix}\right)$, for independent formal parameters $T$ of size $d\times(h-d)$.  Modulo $p$, the column matrix $\binom A C$ has rank $d$.  Choose $T$ so that $A+TC$ is invertible: $A$ is injective on $\ker C$, and one prescribes $T$ on $\im C$ so that the images of a complementary subspace fill a complement to $A(\ker C)$.  Hence $\det(A+TC)$ is a nonzero polynomial and is nonzero at the generic point of the formal parameter space.  In this \emph{covariant} convention, the maximal stable rank of $F$ modulo $p$ is the dimension $d$ of the group.  That maximum is attained here, equivalently the group is ordinary.  Adding back the fixed \etale part and applying Serre--Tate gives an ordinary generic deformation of $A_0$.  If $A_0$ were nonordinary this would contradict the constant stable ranks forced by the stratification.  All arguments may be checked after the faithfully flat extension of the residue field.

For an ordinary deformation, $F$ kills $\omega^{(p)}$, and its induced map on $(H/\omega)^{(p)}$ is invertible.  Thus $\im F$ is a complement to $\omega$.  The two summands are the vector bundles associated with the appropriate \etale $\Fp$-local systems of $A[p]$ and $A^*[p]$, with their invertible $F$- or $V$-structures.  Nilpotent invariance of the \etale site supplies their canonical ordinary stratifications.  They are functorial and horizontal for both operators.

Any other horizontal stratification preserves $\im F$ and also $\im V=\omega^{(p)}$.  Relative Frobenius on the formally smooth universal base, and on its double product, is faithfully flat; therefore preservation of $\omega^{(p)}$ implies preservation of $\omega$.  On either summand, an automorphism comparing two such stratifications is a matrix $G$ with $G\equiv1$ modulo the diagonal ideal $J$ and $G=G^{(p)}$ in a basis of its \etale $\Fp$-structure.  Iteration gives $G\equiv1\pmod{J^{[p^n]}}$ for every $n$, hence $G=1$.  Pullback from the universal base proves uniqueness for every infinitesimal deformation.
\end{proof}

\begin{proposition}[R8, scans 9:5--8]\label{prop:DE-R8}
\status{Proved under stated hypotheses}
Let $k$ be perfect, $W=W(k)$, $W_m=W/p^m$, and let
\[
 \mathcal D=W[F,V],\qquad F\lambda=\sigma(\lambda)F,\quad
 \lambda V=V\sigma(\lambda),\quad FV=VF=p.
\]
For $\mathcal D_n=\mathcal D/(F^{n+1},V^{n+1})$, the constant-term pairing gives an isomorphism of $\mathcal D$-bimodules
\[
 \mathcal D_n\xrightarrow{\sim}
 \Hom_{W_{n+1}}({}_{W_{n+1}}\mathcal D_n,W_{n+1}),\qquad
 y\longmapsto[x\longmapsto\phi(xy)].
\]
The coefficient ring $W_n$ on the handwritten page must be replaced by $W_{n+1}$.  With this correction, the two dualizing functors on modules killed by powers of both $F$ and $V$ agree.
\end{proposition}
\begin{proof}
The normal form is
\[
 \mathcal D_n=W_{n+1}\oplus
 \bigoplus_{i=1}^nW_{n+1-i}F^i\oplus
 \bigoplus_{i=1}^nV^iW_{n+1-i}.
\]
In particular its $W$-length is $(n+1)^2$ and $p^n1\ne0$, which already rules out the proposed $W_n$-linear formula.  The constant-term map is left and right $W$-linear and satisfies
\[
 \phi(Fz)=\sigma\phi(zF),\qquad
 \phi(zV)=\sigma\phi(Vz).
\]
The constant summand pairs with itself by multiplication.  The $F^i$ and $V^i$ summands pair only with each other, through multiplication by $p^i$ and a bijective $\sigma^i$ twist.  They pair perfectly because
$\Hom_{W_{n+1}}(W_{n+1-i},W_{n+1})=p^iW_{n+1}$.

Here are the actions, to remove the two transcription slips and the possible ambiguity of the word ``bimodule.''  For a left $W$-linear functional $u$, set
\[
 (a u)(x)=u(xa),\quad
 (u\lambda)(x)=u(\lambda x),\quad
 (uF)(x)=\sigma^{-1}u(Fx),\quad
 (uV)(x)=\sigma u(Vx).
\]
They satisfy the Dieudonn\'e relations and make the pairing a bimodule map.  Thus the right $V$ action uses $V$, not $F$, and the left $V$ action is precomposition with right multiplication by $V$.

The injections $\mathcal D_n\to\mathcal D_{n+1}$ given by multiplication by $p$ are well defined, since $pF^{n+1}=VF^{n+2}$ and $pV^{n+1}=FV^{n+2}$ in the latter quotient.  They identify the annihilator of $F^r,V^r$ in $I=\varinjlim_n\mathcal D_n$ with $\mathcal D_{r-1}$.  Therefore, if $F^rM=V^rM=0$, adjunction and the pairing identify
\[
 \Hom_{\mathcal D}(M,I)
 =\Hom_{\mathcal D}(M,\mathcal D_{r-1})
 \simeq\Hom_{W_r}(M,W_r).
\]
The anti-involution fixing $W$ and exchanging $F,V$ converts the right module into the usual left Dieudonn\'e dual.  This independently verifies the supplied argument, including its indices.
\end{proof}

\begin{proposition}[R9, scans 6:22--27]\label{prop:DE-R9}
\status{Proved under stated hypotheses}
No ordinary stratification of $H^1_{\dR}$ can be functorial for all elliptic schemes in characteristic $p$ and commute with all base changes.  The three corresponding assertions on independence of a lift, recovery of the classical Dieudonn\'e module, and classification of lifts by their Hodge filtration hold on nilpotent divided-power thickenings, with $p$ locally nilpotent.
\end{proposition}
\begin{proof}
For an explicit obstruction take $k=\overline{\mathbb F}_3$ and the Legendre family
\[
 y^2=x(x-1)(x-\lambda),\qquad \lambda=2+t,
\]
on the formal neighbourhood of $t=0$.  The Hasse polynomial is $1+\lambda=t$.  More generally, for every odd prime, put $m=(p-1)/2$.  The Legendre Hasse polynomial, up to a nonzero scalar depending on the differential convention, is
\[
 H_p(\lambda)=\sum_{i=0}^m\binom mi^2\lambda^i.
\]
It has positive degree, $H_p(0)=1$, and
$H_p(1)=\binom{p-1}{m}\ne0\pmod p$.  Over an algebraic closure choose a root $\lambda_0$, necessarily different from $0,1$, and set $\lambda=\lambda_0+t$.  The discriminant $16\lambda^2(1-\lambda)^2$ is a unit at $t=0$, while the generic Hasse invariant is nonzero.  For $p=2$, use instead
\[
 y^2+txy+y=x^3;
 \qquad\Delta=t^3+1,\qquad\operatorname{Hasse}=t.
\]
These give a smooth family in every characteristic with supersingular closed fibre and ordinary generic fibre.  Thus $F^2=0$ on the closed-fibre de Rham cohomology and $F^2\ne0$ generically.  Functoriality for relative Frobenius would make $F$ horizontal.  Pullback of an ordinary formal stratification by $(0,t)$ would then propagate $F^2=0$ to the whole formal neighbourhood, a contradiction.  This is the argument of the manuscript, also published in \cite[\S3.5]{Grothendieck1968Crystals}, with fully specified families.  The retained program \texttt{code/R7\_R10\_checks.py} verifies the two displayed discriminants and the Legendre coefficient formula by exact symbolic arithmetic; the general binomial identity above supplies the proof for all odd primes.

For the positive assertion write $D(G)$ for the \emph{contravariant} Dieudonn\'e crystal, so that
\[
 0\longrightarrow\omega_G\longrightarrow D(G)_S
       \longrightarrow\operatorname{Lie}(G^*)\longrightarrow0.
\]
For $G=A[p^\infty]$, its evaluation on a PD thickening is canonically $H^1_{\dR}$ of any abelian lift, with the displayed Hodge filtration.  The published locators are \cite[\S2.2, Theorem~2; \S3, Proposition~2 and (3.1)]{BerthelotMessing1979}.  In the original universal-extension construction, see \cite[Chapter~I, Proposition~4.1.7; Chapter~II, (9.2), (13.5), Theorem~15.3]{MazurMessing1974}; the last theorem identifies evaluation over a perfect field with their classical contravariant module.  If instead one uses $M_k(G)=\Hom(G,CW)$ with the precise convention of \cite[\S4.3, Theorem~4 and Remark~1]{BerthelotMessing1979}, that comparison is $\sigma$-semilinear: its $W(k)$-linear form is $\sigma^*M_k(G)\simeq\Gamma D(G)$.  This harmless Frobenius twist must be recorded rather than silently identifying the two conventions.

The lifting equivalence itself is \cite[Chapter~V, Theorem~1.6 for Barsotti--Tate groups; Theorem~1.10 for abelian schemes]{Messing1972}: admissible filtrations are locally direct summands with locally free quotient lifting the special-fibre filtration.  In Messing's original covariant convention use $\operatorname{Lie}E(G)$, or apply it to $G^*$ to obtain the contravariant convention above.  The local-liftability qualification for Barsotti--Tate groups in that book is supplied by \cite[Th\'eor\`eme~4.4 and Corollary~4.7]{Illusie1985BT}; for abelian schemes it is already \cite[Chapter~IV, Lemma~2.8.1]{Messing1972}.

Here ``nilpotent divided powers'' is a condition on the divided-power filtration, not merely nilpotence of the underlying ideal.  The distinction explains why the negative ordinary-stratification argument does not contradict the positive crystalline statement: in characteristic $p$, a PD element $t$ satisfies $t^p=p!\gamma_p(t)=0$, whereas the ordinary formal neighbourhood permits arbitrary powers of $t$.
\end{proof}

\begin{proposition}[R10, scans 7:16--17]\label{prop:DE-R10}
\status{Proved under stated hypotheses}
If $X/k$ is smooth and proper and $k$ is algebraically closed of characteristic $p$, then
\[
 H^i_{\mathrm{str}}(X/k,\OO)
 \simeq H^i(X,\OO_X)_{\ss}
 \simeq H^i_{\et}(X,\Fp)\otimes_{\Fp}k.
\]
For $W=W(k)$, the no-PD infinitesimal cohomology $H^i((X/W)_{\inf},\OO)$ vanishes for $i>\dim X$.  Thus the ordinary abelian example in R10 is valid; the vanishing does not require ordinarity.
\end{proposition}
\begin{proof}
For smooth $X$, formal local lifting identifies ordinary stratifications and infinitesimal crystals.  Frobenius-kernel diagonals are cofinal in the ordinary infinitesimal diagonals, and Frobenius on $X$ is faithfully flat.  Their descent description computes the stable part of proper coherent cohomology.  On a finite-dimensional $k$-space with semilinear Frobenius, the stable image is the scalar extension of its $F=1$ subspace; $F-1$ is surjective.  Artin--Schreier consequently identifies this space with $H^i_{\et}(X,\Fp)\otimes k$.  This is also the $W_1$ case of \cite[Corollary~4.7]{Ogus1975} and of Theorem~\ref{thm:DE-inf}, whose compatible local-embedding argument removes the global-embedding hypothesis.

For $i>\dim X$, coherent cohomology is zero, so the mod-$p$ \etale group is zero.  Induction on $n$ using the coefficient exact sequences gives $H^i_{\et}(X,\mathbb Z/p^n)=0$.  Passing to the inverse limit gives the assertion by Theorem~\ref{thm:DE-inf}; its Mittag--Leffler verification handles the preceding cohomological degree.  This is integral, including torsion, rather than a rational slope-zero assertion alone.
\end{proof}


\subsection{F.2. Source access and literature searches}
Every bibliographic entry below records its actual access state and the
precise checked passages. An opened metadata page is not described as
an opened theorem. The accompanying \path{notes/SOURCE_CENSUS.tsv}
and \path{notes/SOURCE_CENSUS.md} give one row per named or used source,
including unsuccessful attempts. The report incorporates that census.
The three agent search logs and the dated root search logs
\path{notes/ROOT_SEARCH_LOG_04.md}, \path{notes/ROOT_SEARCH_LOG_05.md},
\path{notes/ROOT_SEARCH_LOG_06.md}, \path{notes/ROOT_SEARCH_LOG_07.md},
\path{notes/ROOT_SEARCH_LOG_08.md}, \path{notes/ROOT_SEARCH_LOG_09.md}, and
\path{notes/ROOT_SEARCH_LOG_10.md}, \path{notes/ROOT_SEARCH_LOG_11.md}, and
\path{notes/ROOT_SEARCH_LOG_12.md}, and
\path{notes/ROOT_SEARCH_LOG_13.md}, and
\path{notes/ROOT_SEARCH_LOG_14.md}, and
\path{notes/ROOT_SEARCH_LOG_15.md}, and
\path{notes/ROOT_SEARCH_LOG_16.md}, and
\path{notes/ROOT_SEARCH_LOG_17.md}, and
\path{notes/ROOT_SEARCH_LOG_18.md}, and
\path{notes/ROOT_SEARCH_LOG_19.md}, and
\path{notes/ROOT_SEARCH_LOG_20.md}, and
\path{notes/ROOT_SEARCH_LOG_21.md}, and
\path{notes/ROOT_SEARCH_LOG_22.md}, and
\path{notes/ROOT_SEARCH_LOG_23.md}, and
\path{notes/ROOT_SEARCH_LOG_24.md}, and
\path{notes/ROOT_SEARCH_LOG_25.md}, and
\path{notes/ROOT_SEARCH_LOG_26.md}, and
\path{notes/ROOT_SEARCH_LOG_27.md}, and
\path{notes/ROOT_SEARCH_LOG_28.md}
\path{notes/ROOT_SEARCH_LOG_29.md}
\path{notes/ROOT_SEARCH_LOG_30.md}
record the searches for the requested topics. We make no priority claim.
The source list includes unopened originals to make the limits visible;
their presence is not an assertion that their theorems have been checked.
